\section{Prescribed driving and response diagnostics}\label{app:driving}
\paragraph{Linear coupling to the localized mode.}
For radiation of the same scalar field, a conditional projection identity
makes the distinction between scattering and local forcing precise.
In a fixed real $l=1$ harmonic, complexify the doubled reduced perturbation
$\psi=(rz_l,r\overline z_l)^T$ and write an additive source as $S$.
Equation~\eqref{nl:linear} and its conjugate give
\begin{equation}
 \psi_{tt}+2i\omega\sigma_3\psi_t+\mathcal K\psi=S,
 \qquad
 \mathcal K=-\partial_r^2+\frac{2}{r^2}
 +(D-\omega^2)I+C\sigma_1,
 \label{env:time}
\end{equation}
where $D,C$ are those of Eqs.~\eqref{eq:radialA}--\eqref{eq:radialB},
$\sigma_3=\operatorname{diag}(1,-1)$, and
$\sigma_1=\left(\begin{smallmatrix}0&1\\1&0\end{smallmatrix}\right)$.
The physical perturbation is recovered by the conjugate frequency pairing.
For a physical additive source in this real harmonic block,
$S_2=\overline{S_1}$; independent components are used only in the formal
complexification, with the conjugate pairing restored for physical forcing.
Use the radial inner product with measure $dr$ and the regular
self-adjoint domain of $\mathcal K$. With $X=(\psi,\psi_t)^T$, define
\begin{equation}
 \mathcal G=\begin{pmatrix}0&I\\-\mathcal K&-2i\omega\sigma_3\end{pmatrix},
 \qquad
 \mathcal J=\begin{pmatrix}2\omega\sigma_3&-iI\\iI&0\end{pmatrix}.
 \label{env:generator}
\end{equation}
Then $\mathcal G^\dagger\mathcal J+\mathcal J\mathcal G=0$ on that domain.
For a nonzero localized fundamental mode
$H(\rho)u=0$, where $H(\rho)=\mathcal K-\rho^2I-2\omega\rho\sigma_3$,
put $X_u=(u,i\rho u)^T$. If $\rho>|\omega|$, its conserved form is positive:
\begin{equation}
 N=\langle X_u,\mathcal JX_u\rangle
   =-\langle u,H'(\rho)u\rangle
   =2\langle u,(\rho I+\omega\sigma_3)u\rangle>0.
 \label{env:norm}
\end{equation}
Here $\rho$ is the fundamental frequency, not $2\rho$.
Positivity on this mode is neither a linear nor a nonlinear stability
theorem for the full system.
The corresponding left projection $c=\langle X_u,\mathcal JX\rangle/N$
obeys
\begin{equation}
 \dot c=i\rho c-\frac{i}{N}\langle u,S\rangle,
 \label{env:projection}
\end{equation}
provided the radial Green boundary form vanishes. On $[0,R]$ the
additional term is $i[u'^\dagger\psi-u^\dagger\psi']_0^R/N$.
At the origin this follows from regular values $u,\psi=O(r^2)$ and
regular derivatives $u',\psi'=O(r)$. At infinity the localized mode and
its derivative must decay against the controlled scattering field and
its derivative; this includes suitable finite wave packets.

Thus a homogeneous linear solution ($S=0$) with $c(0)=0$ cannot populate
this projection on the fixed background. Spatial separation of an
incoming packet alone does not imply exact orthogonality, because the
localized mode has tails. A prescribed local additive source can instead
have $\langle u,S\rangle\ne0$ and drive the mode. For $c(0)=0$ and a
finite square-integrable projected source, Eq.~\eqref{env:projection} gives
\begin{equation}
 N|c(T)|^2\leq\frac{T}{N}\int_0^T|\langle u,S(t)\rangle|^2\,dt.
 \label{env:pulse}
\end{equation}
This is a mode-projection bound, not a statement about the energy supplied
by an unspecified environment. Although Eq.~\eqref{env:norm} is positive
on this eigenspace, $\mathcal J$ is indefinite on the full phase space.
No globally passive port model, loading rate, nonlinear formation
mechanism, or indefinite maintenance follows. The local source and its
physical coupling must be specified separately; the nonlinear radiation
obstruction in Section~\ref{nl:section} addresses a different question.

% Result, detuning lemma and text independently reviewed; v0.10 stage.
\paragraph{A prescribed local pulse.}
For the exact T2 mode, write $u=(A,B)^T$ with the reduced real components
of Eq.~\eqref{eq:ansatz}, in the exact-mode Jost normalization specified in
Section~\ref{nl:conditional-monopole}. The projection identity gives a
concrete excitation test for this already existing bound mode.
In dimensionless mass-one units, prescribe a source with real angular
profile $Y_{10}$. Choose the radial envelope
$\chi(r)=(r-1)^2(2-r)^2$ on $1\le r\le2$
and zero elsewhere, and temporal envelope $\eta(t)=\sin^2(\pi t)$ on
$0\le t\le1$ and zero otherwise. The physical source in the original field equation is
$J_\phi=\epsilon\chi Y_{10}\eta\exp[i(\omega+\rho)t]$, with real
$\epsilon$, the small strength of the external drive in this linear-response
calculation. It is added to the right-hand side of Eq.~\eqref{nl:model}.
The reduced source in Eq.~\eqref{env:time} is
$S=\epsilon r\chi\eta(e^{i\rho t},e^{-i\rho t})^T$;
its conjugate channel is not independent and must be retained. For the fixed real
radial mode components $A,B$, define
\[
 I_A=\int_1^2 r\chi(r)A(r)\,dr,\qquad
 I_B=\int_1^2 r\chi(r)B(r)\,dr.
\]
The resulting finite-pulse overlap is
\[
L(0)=\frac{I_A}{2}+I_B\int_0^1\eta(t)e^{-2i\rho t}\,dt.
\]
With $N>0$ from Eq.~\eqref{env:norm}, vanishing initial projection and
vanishing Green boundary form, the pulse leaves
$c(1)=-i\epsilon e^{i\rho}L(0)/N$.
Exact integration of the stored nominal polynomials, combined with their
certified componentwise error bounds and the same certified frequency box,
gives an enclosure separated from zero. Conditional on the same exact T2
mode identification and normalization,
a conservative bound is $|L(0)|>1/500$. This certifies nonzero linear mode
projection for nonzero $\epsilon$, vanishing initial projection and the
stated boundary conditions; it does not determine an absolute excitation
energy without the mode norm. Replacing $Y_{10}$ by $Y_{00}$ gives zero
direct dipole projection by angular orthogonality; other modes are not excluded.

% Detuning bound independently checked in DETUNING-REVIEW.txt.
\paragraph{Finite detuning.}
The conclusion need not require an exactly tuned external source. Replacing
its carrier by $\omega+\rho+\delta$ gives
\[
 L(\delta)=I_A\int_0^1\eta(t)e^{i\delta t}\,dt
 +I_B\int_0^1\eta(t)e^{-i(2\rho+\delta)t}\,dt.
\]
Since $\int_0^1t\eta(t)\,dt=1/4$,
\[
 |L(\delta)-L(0)|\le\tfrac14(|I_A|+|I_B|)|\delta|.
\]
The conservative bounds $|I_A|<3/1000$, $|I_B|<6/1000$ therefore imply
$|L(\delta)|>7/8000$ for $|\delta|\le1/2$. This is a sufficient excitation
bound for the prescribed local source, not a linewidth or a decay rate.
Its duration is dimensionless, not a measured physical loading time.
Neither calculation describes formation of the underlying Q-ball,
self-sustained excitation, or coupling to a specified cosmological field.

\paragraph{Canceling one excitation is not stopping the Q-ball.}
The source-projection identity also supplies a simple linear control prediction.
Let $d=e^{-i\rho t}c$, initially zero, and apply two identical physical pulses
of the preceding family, with the same local rotating-frame waveform and
strength $\epsilon$, separated by $\Delta\geq1$.
Their total projected drive is
\begin{equation}
 I_{\mathrm{pair}}=\epsilon L(\delta)
                 \bigl(1+e^{-i\rho\Delta}\bigr).
 \label{env:counterpulse}
\end{equation}
Consequently $\Delta=\pi/\rho>1$ cancels the mode projection generated by the
first pulse, even for a common local detuning $\delta$. Both sideband terms
in $L$ are retained. This construction needs identical pulses in the
rotating frame: in laboratory variables their sources obey
$J_2(t)=e^{i\omega\Delta}J_1(t-\Delta)$.
A sign-reversed second pulse at that same half-period instead doubles the
rotating projected drive. For $|\delta|\leq1/2$, the preceding bound ensures
a nonzero first-pulse projection. Relative to that projection, a timing
error $h$ leaves the remainder $2|\sin(\rho h/2)|\leq\rho|h|$.
These statements are analytic consequences of the fixed linear model.
The following finite-time calculation tests the zero-detuning case only;
it does not numerically establish the entire detuning or timing-error family.
The common unknown modal norm cancels
from the cancellation condition; extinguishing an arbitrary pre-existing
excitation would require its complex amplitude and calibration.
Other modes, continuum waves, radiation already emitted, and nonlinear
backreaction are not canceled by this one condition. Neither zero modal
projection nor zero linear radiation implies that the charged Q-ball itself
has stopped rotating or disappeared.

% Candidate only. No canonical paper edit; figure awaits rendering/review.
\paragraph{A finite-time radial counterpulse test.}
We tested the two-pulse prediction by evolving the linearized complex radial
$\ell=1$ equation on the fixed nominal T2 background, with the conjugate
physical channel retained.
Four initially zero arms used no pulse, one pulse, two equal-sign local
rotating-frame pulses separated by $\Delta=\pi/\rho$, and a sign-reversed
second pulse at the same delay.
The calculation used $R=20$, $T=6$, grid spacings $h=0.04,0.02$ and
fourth-order Runge--Kutta time steps $h/8$, with forcing reported per unit
linear drive amplitude. It is a radial linear PDE calculation, not a
nonlinear three-dimensional evolution.

Both grids passed the predeclared projection-balance and comparison checks.
The equal-sign pair reduced the final target projection below $2\%$ of
the single-pulse reference, while the sign-reversed pair agreed with twice
the reference complex projection within $2\%$.
These are numerical acceptance bounds, not interval-certified errors.
The observed fine-grid ratio was approximately $3.61\times10^{-5}$
($0.0036\%$); the coarse-grid ratio was approximately $1.44\times10^{-4}$.
The balance includes the explicitly computed defect of the stored nominal
mode on the spatial grid; the source-only discrepancy is reported separately.
The global second-component defect norm, which includes the forced cutoff
of the nominal mode at the finite boundary, actually increases on refinement;
we do not claim
global eigenfunction-error convergence. The short-time paired defect and
projection checks concern the evolving, spatially confined response.
Figure~\ref{fig:counterpulse-dynamics} also shows that the radial field
norm need not vanish when the selected projection is suppressed.
At $T=6$ its fine-grid value was $0.0119332$ for the equal-sign pair,
compared with $0.00818567$ for the single pulse.
The mode-subtracted endpoint norm uses the nominal reconstruction
$q-cA-\overline cB$ and is not an orthogonal energy decomposition.
No total energy, charge removal, long-time stability or complete stillness
of the charged background is inferred.

\begin{figure}[htbp]\centering
\includegraphics[width=.98\linewidth]{counterpulse-dynamics.pdf}
\caption{Stored radial linear-PDE histories for the T2 counterpulse test.
Top: target projection magnitude, normalized by the fine-grid final
single-pulse value. Bottom: reduced-field $L^2(dr)$ norms; endpoint symbols
give norms after nominal mode subtraction. Colors identify pulse arms,
solid/dashed lines the two grids. The exactly zero null arm is omitted
from the logarithmic panel. Shading marks pulse support; only paired arms
receive the second pulse. Lines connect recorded times without a new
evolution. The laboratory second pulse includes the carrier-phase correction
of the analytic construction. A small target projection coexists with a
remaining field; these norms are neither radiated energy nor Noether charge.}
\label{fig:counterpulse-dynamics}
\end{figure}

% Candidate for v0.14; not a canonical edit or publication approval.
% RESULT.json SHA256 7c9ff31f20b4ce7ee65e146a13ba1ccdda18ae04c07757e7ad1abf56b4702e88
% RESULT-REVIEW.txt SHA256 4bb05fb4bcaa8c187c7fd850c0d2908283bedf98b3e1d494e5e3676658373e33
\paragraph{Changing the complete pulse shape.}
A subsequent fixed comparison used the same nominal linear radial
$\ell=1$ evolution, $R=20$, $T=6$ and $h=0.04,0.02$.
It compared the annular drive above with a mode-adapted drive, using
single pulses and equal-sign pairs separated by $\pi/\rho$.
The new drive retains both sideband terms of a quintic turn-on construction
and has a spatial cutoff equal to one for $r\leq4$ and zero for $r\geq8$.
The reference instead uses the annulus $1<r<2$ and a $\sin^2$ temporal
window with one rotating frequency. Thus both spatial and temporal
waveforms change; the comparison does not isolate a spatial cause.

Each single pulse was normalized to the same reduced-source norm
\begin{equation}
 B_{\mathrm{src}}=\int\!\int |s(r,t)|^2\,\dd r\,\dd t,
 \qquad B_{\mathrm{src}}\simeq1.35281\times10^{-3}
 \quad(h=0.02).
\end{equation}
Pairs have twice this budget. The integrated budgets agreed with their
prescribed values within $6\times10^{-9}$ relatively on both grids.
This norm is neither mechanical work nor injected energy or charge,
and equal norms do not imply equal actuator support or hardware cost.
For the single pulse, the fine-grid diagnostic $|d(T)|^2/B_{\mathrm{src}}$
is about $32$ times larger for the shaped drive
(Table~\ref{tab:source-shaping}); this is not an energy-efficiency ratio
or an optimality result.

\begin{table}[htbp]\centering\small
\caption{Fixed source-shaping comparison at $T=6$, $h=0.02$.
Here $p=\partial_tq$, and ``remainder'' means subtraction of the nominal
mode reconstruction. Field norms are $L^2(\dd r)$ norms, not energies;
the pair rows give the total remaining field and velocity norms.}
\label{tab:source-shaping}
\begin{tabular}{lrr}\toprule
Diagnostic&Annular drive&Shaped drive\\\midrule
Single-pulse $|d|^2/B_{\mathrm{src}}$&$0.001214$&$0.03872$\\
Single-pulse $q$ remainder norm&$0.007992$&$0.002596$\\
Single-pulse $p$ remainder norm&$0.01502$&$0.004954$\\
Pair $\|q\|$&$0.01193$&$0.001055$\\
Pair $\|p\|$&$0.02153$&$0.001930$\\\bottomrule
\end{tabular}
\end{table}

The shaped pair leaves approximately $8.8\%$ and $9.0\%$ of the
reference pair's field and velocity norms, respectively; both remain
nonzero. Nominal mode subtraction is not an orthogonal energy
partition. Both grids passed the fixed comparison and defect-balance
checks, but the global second-component nominal eigenmode residual norm increases
from approximately $0.529$ to $1.495$ on refinement. The successful
paired-defect balance therefore does not establish global eigenfunction
convergence. These finite-box, short-time results show reduced residual
motion for the specified composite drive; they establish neither full
stillness of the field nor removal of the charged background, nonlinear
stability or a lifetime. This original two-source comparison did not
separate changes of spatial and temporal design. The subsequently
specified cross-comparison below addresses that limitation within a
fixed actuator family.

% Reviewed stored-profile attribution, incorporated in draft 0.17.
% RESULT 03e1c3129984f7b96713f901e7e7767b5f32d4bb6b33afa380178d34dc8ffbb7
% RESULT-REVIEW 3b665e564d1471fcfbf80cb42d118615c011bfb7b4a2cffea75522ca27619b9c
A subsequent stored-profile audit identifies the origin of the growing
second-component defect in this box. The exported reference mode has
$v(20)=0.004228503237$, whereas the evolution operator imposes a zero
Dirichlet value. For the three-point stencil this discrepancy adds
$b_j=\delta_{j,n-1}v(20)/h^2$ to the continued-ghost residual, giving
$\|b\|_h^2=|v(20)|^2/h^3$. Its squared norm is $0.2793787442$ and
$2.2350299533$ on the two grids; the continued residual squared norms
are $1.52\times10^{-8}$ and $9.50\times10^{-10}$, with signed cross
terms $-9.74\times10^{-10}$ and $-4.84\times10^{-10}$.
The boundary discrepancy therefore dominates this component's growing
residual. It does not dominate the first component. This is an algebraic
attribution, not a modification of the evolution boundary condition or a
new eigenfunction calculation. The original global defect remains valid
for the original operator. Two decreasing interior residuals neither
certify the exported mode nor bound the integrated influence of the
boundary discrepancy on the pulse response.

\begin{figure}[htbp]\centering
\includegraphics[width=.98\linewidth]{source-shaping.pdf}
\caption{Endpoint diagnostics from the fixed spatiotemporal pulse comparison.
Both spatial support and temporal waveform change. The common source norm
is not energy or work. The shaped source increases the selected projection
per source norm and reduces the residual response in this comparison;
the complete field does not vanish. Bars show the stored fine-grid result,
not a new evolution or a certified continuum limit.}
\label{fig:source-shaping}
\end{figure}

% Factorial RESULT477978704a19d5b2f2b474c41cf5917367335c521902e0f712303f7a563ec263
% Non-author RESULT-REVIEW66fe64890bbbc3568c19a33610968cd7360f68462ea8f7fb7b467b900dd02a28
\paragraph{A fixed factorial actuator comparison.}
The new comparison uses two spatial pairs, $A=(w,0)$ and
$B=(\zeta u,\zeta v)$, where $w$ is the preceding annulus profile,
$\zeta$ the cutoff defined above (one below $r=4$ and zero above $r=8$),
and $u,v$ the stored nominal sidebands.
Two temporal pairs are prescribed on $0<t<1$:
\begin{align}
 T_0&=\bigl(\sin^2(\pi t)e^{i\rho t},
                  \sin^2(\pi t)e^{-i\rho t}\bigr),\\
 T_1&=\bigl([g''+2i(\rho+\omega)g']e^{i\rho t},
                  [g''-2i(\rho-\omega)g']e^{-i\rho t}\bigr),
 \qquad g=10t^3-15t^4+6t^5.
\end{align}
Both temporal pairs vanish outside $0<t<1$; $g$ is extended constantly
as zero before the pulse and one afterward.
Each complex source is the component sum
$s_{ij}=\mathcal N_{ij}(R_{i1}T_{j1}+R_{i2}T_{j2})$,
where $R_i=A$ or $B$. There is no conjugation in this sum; the second
physical equation receives $s_{ij}^*$.
The positive factor $\mathcal N_{ij}$ fixes the same single-pulse
$B_{\mathrm{src}}$ \emph{after} summation, retaining component interference.
The two old corners $A/T_0$ and $B/T_1$ are implementation controls;
$A/T_1$ and $B/T_0$ are the new crossed drives. Each uses both a single
pulse and an equal-sign pair at the unchanged delay $\pi/\rho$.
The nominal operator, $R=20$, $T=6$, $h=0.04,0.02$ and time step $h/8$
remain fixed. Both grids passed the predeclared numerical checks.

\begin{table}[htbp]\centering\small
\caption{Four fixed normalized actuator designs at $T=6$, $h=0.02$.
The projection diagnostic is for a single pulse; the final two columns
are total remaining pair norms, not mode-subtracted norms or energies.}
\label{tab:source-factorial}
\begin{tabular}{lrrr}\toprule
Design&$|d_{\mathrm{single}}|^2/B_{\mathrm{src}}$&Pair $\|q\|$&Pair $\|p\|$\\\midrule
$A/T_0$&$0.001214$&$0.01193$&$0.02153$\\
$A/T_1$&$0.0004100$&$0.006390$&$0.01426$\\
$B/T_0$&$0.1586$&$0.02024$&$0.005692$\\
$B/T_1$&$0.03872$&$0.001055$&$0.001930$\\\bottomrule
\end{tabular}
\end{table}

Within this family, $B/T_0$ gives the largest single-pulse projection
diagnostic, whereas $B/T_1$ leaves the smallest pair norms in both $q$
and $p$. Strong projected excitation and small residual field are thus
different objectives. The spatial-pair change $A\to B$ increases the
pair $q$ norm at $T_0$ but decreases it at $T_1$: its effect depends on
the temporal partner. These are conditional design comparisons, not
universal additive spatial and temporal effects. In particular $A$
omits a spatial component, $T_1$ changes sideband weights and phase
structure, and normalization depends on both factors. Equal source norm
still implies neither equal work nor energy efficiency. All four pairs
suppress their selected projection under the fixed criterion, while their
full responses differ and remain nonzero. The nominal-mode defect and
finite-box limitations above remain. That fixed-phase comparison did not
test phase-error robustness; it established neither exact preparation,
complete stillness nor a lifetime.

\begin{figure}[htbp]\centering
\includegraphics[width=.98\linewidth]{factorial-comparison.pdf}
\caption{Stored endpoint diagnostics for the fixed $2\times2$ actuator
family. The largest single-pulse projection per source norm occurs at
$B/T_0$, while $B/T_1$ gives the smallest remaining pair field and
velocity norms. The source is normalized after adding its two components;
these comparisons are not energy efficiencies or a universal separation
of independent spatial and temporal mechanisms. No additional evolution
or phase-error test is represented by this figure.}
\label{fig:source-factorial}
\end{figure}

% Phase RESULTf492a34c52a608a36d3d6054c1342cc6bf3dd963021e3576d5d6447423a25301
% Reviewed ea5b091239d228e8da5ac68c60108c238c3052a2ee4272f9012366e61459ec37
% Endpoint SUMMARY3329b5213520fe157b71c0a79d530a36e03ce9a644fd13d48cc71031e0c53aae
\paragraph{A constant relative pulse phase.}
A further bounded test changes only the additional phase $\theta$ of the
second pulse, keeping the delay and the carrier-phase correction fixed.
The coupling to $q^*$ makes the nominal Bogoliubov evolution real-linear,
not complex-linear: the response to $i s$ need not equal $i q$.
For each of the four fixed designs, one new delayed-$i s$ response
therefore completes the real basis. At the same endpoint,
\begin{equation}
 q_\theta=q_{\mathrm{single}}
   +\cos\theta\,(q_{\mathrm{pair},0}-q_{\mathrm{single}})
   +\sin\theta\,q_{\mathrm{delayed},90}.
 \label{eq:phase-real-basis}
\end{equation}
The same identity holds for $p$ and the real-linear projection $d$.
Four directly evolved $\theta=\pi/4$ pairs check this reconstruction:
the largest reported relative discrepancy in $q,p,d$ across both grids
was $1.14\times10^{-13}$. This tests a known linear identity, not a new
physical effect. Other constant phases below are evaluations of stored
endfields, not additional time evolutions. The setting remains
$R=20$, $T=6$, $h=0.04,0.02$.

\begin{table}[htbp]\centering\small
\caption{Selected constant-phase endpoints for $B/T_1$ at $T=6$,
$h=0.02$. Norms are total remaining radial fields, not energy.
The phase values were fixed before the endpoint analysis.}
\label{tab:phase-selected}
\begin{tabular}{rrr}\toprule
$\theta$&$\|q\|$&$\|p\|$\\\midrule
$0^\circ$&$0.001055$&$0.001930$\\
$-5^\circ$&$0.001062$&$0.002068$\\
$+5^\circ$&$0.001423$&$0.002462$\\
$-15^\circ$&$0.001986$&$0.003720$\\
$+15^\circ$&$0.002562$&$0.004413$\\
$45^\circ$&$0.006301$&$0.01125$\\\bottomrule
\end{tabular}
\end{table}
These values need not be symmetric about zero phase. Cancellation of the
selected coefficient and minimization of the full field norm are different
conditions; no new minimum search was performed.

For a uniformly distributed phase that is \emph{constant within each
pulse pair}, write any endpoint field as $a+b\cos\theta+c\sin\theta$.
Its root mean square norm is exactly
\begin{equation}
 \sqrt{\mathbb E_\theta\|a+b\cos\theta+c\sin\theta\|^2}
 =\sqrt{\|a\|^2+\tfrac12(\|b\|^2+\|c\|^2)}.
\end{equation}
The Gram-matrix evaluation was also checked against eight equally spaced
phases. Table~\ref{tab:phase-rms} keeps all four designs in the comparison.
\begin{table}[htbp]\centering\small
\caption{Uniform constant-phase RMS endpoint norms, $T=6$, $h=0.02$.
These are square roots of mean squared norms, not mean norms.}
\label{tab:phase-rms}
\begin{tabular}{lrr}\toprule
Design&$\sqrt{\mathbb E_\theta\|q\|^2}$&$\sqrt{\mathbb E_\theta\|p\|^2}$\\\midrule
$A/T_0$&$0.01128$&$0.02168$\\
$A/T_1$&$0.007123$&$0.01447$\\
$B/T_0$&$0.02552$&$0.04088$\\
$B/T_1$&$0.01106$&$0.01954$\\\bottomrule
\end{tabular}
\end{table}
Although $B/T_1$ leaves smaller pair norms than $A/T_1$ at zero phase,
its two RMS norms are larger in this ensemble. Its fixed-phase advantage
is thus conditional on phase preparation.
The comparison fixes source norm, not the previously excited target
amplitude: $A/T_1$ has a much smaller single-pulse projection diagnostic
($0.0004100$ versus $0.03872$ for $B/T_1$). Its lower RMS therefore does
not establish a better storage or erasure protocol at equal preparation.
An RMS ranking is not necessarily
an ordering of mean norms. This ensemble is neither a time average nor a
simulation of fluctuating environmental noise. The constant phase does
not change the disjoint pulses' source budget $2B_{\mathrm{src}}$, but that
budget remains distinct from work or energy. The global nominal eigenmode
residual limitation remains unchanged; neither exact BIC preparation,
complete stillness nor nonlinear or long-time robustness follows.

\begin{figure}[htbp]\centering
\includegraphics[width=.98\linewidth]{phase-response.pdf}
\caption{Constant-phase endpoint reconstruction for all four fixed
actuator designs, from their stored real-linear response bases.
The two grids compare nominal finite-box observables; plotted phase
samples are not independent time evolutions. The additional phase is
constant within a pulse pair, not a time-dependent noise source.
Field norms and uniform-phase RMS values do not measure energy cost.}
\label{fig:phase-response}
\end{figure}

\paragraph{Quadratic endpoint energy and the rotating generator.}
A further analysis uses the stored single- and pair-pulse endpoints,
without repeating the time evolution. With $p=q_t$ and the same
$h$-weighted inner product and Dirichlet operator as the evolution,
define
\begin{align}
 H_2&=\|p\|_h^2+\langle q,Kq\rangle_h
       +\operatorname{Re}\,h\sum_j C_j q_j^2,
       \label{eq:endpoint-generator}\\
 Q_2&=2\operatorname{Im}\langle q,p\rangle_h
       +2\omega\|q\|_h^2,\qquad E_2=H_2+\omega Q_2.
       \label{eq:endpoint-energy}
\end{align}
Here $K=-\Delta_h+2/r^2+1-4f^2+(9/2)f^4-\omega^2$ and
$C=-2f^2+3f^4$. There is no factor $1/2$: $H_2$ is the coefficient
of second order in $E-\omega Q$ for the stated action convention,
using one physical complex field rather than counting its conjugate twice.
This rotating quadratic generator need not be positive.

\begin{table}[htbp]\centering
\caption{Stored endpoint coefficients at $T=6$, $h=0.02$.
All entries are in units of $10^{-4}$, rounded to four significant digits.
Within a fixed pulse count the designs have equal source-L2 budgets;
a pair has twice the single-pulse budget.}
\label{tab:endpoint-energy}
\begin{tabular}{lrrrr}\toprule
Design&$H_2$ single&$H_2$ pair&$E_2$ single&$E_2$ pair\\\midrule
$A/T_0$&4.782&9.315&5.958&11.51\\
$A/T_1$&2.668&5.180&3.614&6.910\\
$B/T_0$&8.917&1.382&6.467&8.608\\
$B/T_1$&2.129&0.05203&1.201&0.03759\\\bottomrule
\end{tabular}
\end{table}
The distinction matters even for the direction of change:
$B/T_0$ decreases $H_2$ while increasing $E_2$ from single to pair.
Its $Q_2$ changes from approximately $-2.821\times10^{-4}$ to
$8.319\times10^{-4}$. For $B/T_1$ both endpoint coefficients decrease;
for the two $A$ designs both increase. These directions also occur on
$h=0.04$, but two-grid differences are not validated error bounds.

Equal source-L2 budgets do not imply equal work, and the single-to-pair
comparison itself is not a same-budget comparison. For the fixed
semidiscrete model, $\dot H_2=2\operatorname{Re}\langle p,S\rangle_h$.
Thus $H_2(T)$ at zero initial data predicts net rotating-generator work
by an identity; that endpoint-only analysis did not measure an independent
time-integrated power balance. The subsequent test below adds this observable.
The quadratic laboratory source power additionally contains
$2\omega\operatorname{Im}\langle q,S\rangle_h$, while $\dot E_2$
also contains the background-exchange term
$2\omega\operatorname{Im}\,h\sum_j C_jq_j^2$.
In particular, $E_2$ is a coefficient on the linear ansatz, not the complete
laboratory energy change including a second-order background response.
No positive or gross actuator work, energy efficiency, or full shutdown
of the Q-ball is established by Table~\ref{tab:endpoint-energy}.
\paragraph{Time-integrated source work in the nominal linear model.}
A subsequent bounded run retains the same eight single/pair arms,
$R=20$, $T=6$, $h=0.04,0.02$ and $\Delta t=h/8$, and adds power
accumulators at the four Runge--Kutta stages. With the conventions of
Eqs.~\eqref{eq:endpoint-generator}--\eqref{eq:endpoint-energy}, these
integrate
\begin{align}
 P_H&=2\operatorname{Re}\langle p,S\rangle_h,&
 J_Q&=2\operatorname{Im}\langle q,S\rangle_h,&
 X_Q&=2\operatorname{Im}\,h\sum_jC_jq_j^2,\\
 P_{E,\mathrm{src}}&=P_H+\omega J_Q,&
 P_{E,\mathrm{ex}}&=\omega X_Q.
\end{align}
The independently evaluated state coefficients satisfy the numerical
balances $\Delta H_2=\int P_H\,\dd t$,
$\Delta Q_2=\int(J_Q+X_Q)\,\dd t$, and
$\Delta E_2=\int(P_{E,\mathrm{src}}+P_{E,\mathrm{ex}})\,\dd t$.
Unlike the preceding endpoint-only analysis, these integrals use the
actual time-dependent fields and source at the integration stages.
The exchange term remains active after the source ends.

The first pulse is identical in the single and pair arms; the first and
second pulse windows are disjoint. Their trajectories therefore agree before
the second pulse. Subtracting their independently integrated total source
works gives the second pulse's net work on the already excited field:
\begin{equation}
 W_{H,2}=W_{H,\mathrm{pair}}-W_{H,\mathrm{single}},\qquad
 W_{E,\mathrm{src},2}=W_{E,\mathrm{src},\mathrm{pair}}
                       -W_{E,\mathrm{src},\mathrm{single}}.
\end{equation}
This is not the work of an isolated second pulse applied to a fresh
zero state, nor merely a relabeling of endpoint energy differences.

\begin{table}[htbp]\centering
\caption{Net second-pulse source work inferred from time-integrated
pair-minus-single works at $h=0.02$. Units are $10^{-4}$, as in
Table~\ref{tab:endpoint-energy}; entries are rounded to four significant
digits. Negative values denote net extraction in the specified quadratic
source-work convention, not efficiency of a physical actuator.}
\label{tab:second-pulse-work}
\begin{tabular}{lrr}\toprule
Design&$W_{H,2}$&$W_{E,\mathrm{src},2}$\\\midrule
$A/T_0$&$4.533$&$4.207$\\
$A/T_1$&$2.512$&$3.602$\\
$B/T_0$&$-7.535$&$5.939$\\
$B/T_1$&$-2.077$&$-1.122$\\\bottomrule
\end{tabular}
\end{table}
Thus the second $B/T_1$ pulse is a net source-work sink in both
conventions. In contrast, $B/T_0$ removes rotating-generator work while
its laboratory source-work coefficient is positive. The change in $E_2$
also includes the difference of integrated exchange terms; it is not
identical to $W_{E,\mathrm{src},2}$.
Positive and negative parts of the total spatially integrated power
were separately accumulated, so a small net value need not mean little
positive input or little extraction. This splitting applies after spatial
summation, rather than to positive and negative local power densities.

Both grids pass the predeclared $10^{-6}$ normalized balance tolerance,
and the new endfields agree with the stored reference endpoints within
the fixed $10^{-8}$ single-response-normalized tolerance. The reported
scaled two-grid differences of work contributions are below
$4.4\times10^{-4}$, within the prescribed $0.03$ consistency limit.
These are discrete numerical checks, not rigorous continuum error bounds;
the positive/negative power split also need not inherit a smooth fourth-order
convergence rate. The original nominal-mode boundary defect is unchanged.
Equal source-L2 budgets are still not equal work. Finally, the laboratory
source-work coefficient and $E_2$ belong to the frozen linear ansatz:
second-order background response, full nonlinear laboratory energy,
actuator losses and conversion efficiency are not determined. No complete
shutdown or long-time stability follows from this short-window balance.
\paragraph{Phase dependence of the second-pulse source work.}
For the same four source designs, we vary the constant phase $\theta$ of
the delayed pulse while keeping its source norm fixed. Real linearity of
the frozen evolution gives the quadratic source-work coefficient
\begin{equation}
 W_{E,\mathrm{src},2}(\theta)
 =c_0+c_1\cos\theta+s_1\sin\theta
       +c_2\cos(2\theta)+s_2\sin(2\theta).
 \label{eq:phase-work-five}
\end{equation}
Writing the measured quadrature coefficients as $(A,B,D,E,F)$, these are
$(c_0,c_1,s_1,c_2,s_2)=((D+E)/2,A,B,(D-E)/2,F/2)$.
The five coefficients are integrated from the first-pulse response and two
delayed quadrature responses, with a direct $45^\circ$ two-pulse control,
through $T=3$ after both source windows end ($\Delta+1<3$), using
$h=0.04,0.02$ and $\Delta t=h/8$. This is not a renewed $T=6$ energy history.
Thus the plotted curves are postprocessing of measured coefficients, not
thousands of independent evolution runs.

On both nominal grids, only B/T1 has a negative-work interval:
approximately $-66.53^\circ<\theta<57.03^\circ$.
For the fine grid its second-pulse net source work is
$-1.1224\times10^{-4}$ at $0^\circ$ and
$-3.9278\times10^{-5}$ in the direct $45^\circ$ control.
This broad nominal interval concerns signed net source work; it does not
assert equally small residual fields throughout the interval.
In particular, the previously reported residual field norms increase at
$45^\circ$ even though this work remains negative.
All four uniform \emph{constant-phase} means $c_0$ are positive;
for B/T1 the fine-grid mean is $1.1334\times10^{-4}$.
This phase ensemble is not a model of time-dependent phase noise.

Polynomial sign isolation treats the measured floating-point coefficients
as exact rational dyads. Its exactness applies to those nominal polynomials,
not to their discretization errors or to the continuum or nonlinear field
equations. The result concerns quadratic source work in the frozen linear
ansatz, with background exchange distinguished separately; it is neither a
full nonlinear laboratory-energy balance nor an actuator-efficiency claim.

\begin{figure}[htbp]
\centering
\includegraphics[width=\linewidth]{phase-work.pdf}
\caption{Second-pulse source work versus constant phase. Left: all four
designs; solid curves use $h=0.02$, dashed curves $h=0.04$.
Right: B/T1 with nominal fine-grid zero boundaries (dotted), its uniform
constant-phase mean (dash-dotted), and direct $45^\circ$ controls (symbols).
Work is displayed in units of $10^{-4}$. Shading denotes the nominal
negative-work window, not certified physical tolerances.}
\label{fig:phase-work}
\end{figure}

\clearpage
\paragraph{One prescribed temporal phase drift.}
The constant-phase ensemble above does not determine the response to a
phase varying during a pulse. As a separate finite comparison, we replaced
the same delayed B/T1 source $S_d(t)$ by
$S_{\delta,\theta_0}(t)=\exp[\ii\theta_0+\ii\delta(t-\Delta)]S_d(t)$,
with the single preselected value $\delta=0.1$ against $\delta=0$.
The multiplier acts on the entire complex source; the original envelope,
carrier construction and delay are unchanged. It preserves the source
magnitude pointwise and its integrated $L^2$ budget. For each detuning we
independently evolved the responses to $S_{\delta,0}$ and
$\ii S_{\delta,0}$ through $T=3$. Real linearity gives the response at
$\theta_0$ as the cosine--sine combination of these two responses;
the second is not assumed to be $\ii$ times the first.

Averaging uniformly over the \emph{constant initial} phase $\theta_0$
eliminates both the mixed quadrature term and the work cross term with
any deterministic first-pulse field. Consequently the mean second-pulse
laboratory source work is one half of the sum of the two isolated
quadrature works. No first-pulse evolution was repeated; this reduction
applies to the mean, not to phase-specific two-pulse work. With
$(h,\Delta t)=(0.04,0.005)$ and $(0.02,0.0025)$, both means are positive
on both discretizations. On the finer grid the mean increases from
$1.133373236\times10^{-4}$ to $1.248632464\times10^{-4}$,
a difference of $1.152592277\times10^{-5}$ (about $10.2\%$).

The increase exceeds the predeclared diagnostic comparison guard
$7.94\times10^{-11}$, which combines the change between discretizations,
unshifted-work reproduction errors and energy-balance residuals.
This guard is not a certified error bound; space and time resolution change
together. The result establishes sensitivity of this nominal mean source
work to one prescribed phase drift at equal source norm. It does not give
a phase-noise or thermal model, an efficiency gain, or a claim about all
detunings. The frozen-profile approximation, artificial $R=20$ boundary
and omitted nonlinear backreaction remain unchanged.
\paragraph{Local direction and one mirror detuning.}
A follow-up fixed $d=0.1$ and added only the mirror source at $-d$ and
the parameter tangent at zero. Writing $z=\partial_\delta q|_0$, the
exact differentiated finite-grid equation has the same real-linear
operator and forcing $\ii(t-\Delta)S_d$. Its numerical integration retains
both derivatives of the work integrand, from the response and from the
source, for each independently evolved quadrature. The zero-detuning
basis trajectories are needed for these products; the first pulse is
not repeated, and the $+d$ result is reused. For the uniform initial-phase
mean $m(\delta)$, the finer grid gives $m'(0)\simeq1.1006065\times10^{-4}$
in model units. Both grids resolve a positive slope with the prescribed
empirical guard. The three fine-grid means remain positive:
\begin{center}
\begin{tabular}{rr}
\toprule
$\delta$ & Mean laboratory source work $m(\delta)$\\
\midrule
$-0.1$ & $1.0285198\times10^{-4}$\\
$0$ & $1.1333732\times10^{-4}$\\
$+0.1$ & $1.2486325\times10^{-4}$\\
\bottomrule
\end{tabular}
\end{center}

Define the finite odd and even changes by
$O_d=[m(d)-m(-d)]/2$ and
$E_d=[m(d)+m(-d)-2m(0)]/2$.
Their fine-grid values are approximately $1.1005634\times10^{-5}$ and
$5.20289\times10^{-7}$; the residual
$R_d=O_d-dm'(0)\simeq-4.31\times10^{-10}$ is small but not an exact
identity or a certified Taylor remainder. The resolved positive odd
change agrees with the local direction. This is a directional response
of the fixed nominal actuator, not an identification of rotation as its
cause: a free complex accelerator $q_{tt}=\exp[\ii(\kappa+\delta)t]$ already
admits a nonzero detuning slope. More generally, uniform initial-phase
averaging selects the normal causal response, whose work derivative is
a lag-weighted imaginary response kernel; the carrier and field response
enter it jointly. The two-grid guards are diagnostic, and the inherited
profile, boundary and linearization limitations persist. No thermal,
new-memory-mediator or efficiency conclusion follows.

\paragraph{Prescribed finite phase coherence.}
A fluctuating drive phase can be specified without identifying it with
thermal noise. Let the finite-duration source be
$S_\theta(t)=e^{i\theta(t)}S_d(t)$, where the initial phase is uniform and
independent and $\theta(t)-\theta(\Delta)=\sqrt{2D}\,B_{t-\Delta}$ during
the pulse, with $B$ a standard real Wiener process independent of that initial phase. Then
$\mathbb E e^{i[\theta(t)-\theta(s)]}=e^{-D|t-s|}$ and the anomalous
source correlation vanishes. Every realization retains the source's
$L^2$ budget. This does not fix the work delivered to the field.

For the fixed finite real-linear response used above, let $m(\delta)$ denote
the phase-averaged source work when the \emph{entire} prescribed source
is multiplied by $e^{i\delta(t-\Delta)}$. Its phase-diffusion mean obeys
\[
 m_D=\int_{-\infty}^{\infty}
 \frac{D}{\pi(D^2+\delta^2)}\,m(\delta)\,d\delta,
 \qquad D>0.
\]
This follows from the characteristic function of the Cauchy distribution
and the quadratic response kernel. At finite time the finite-dimensional
propagator is bounded; compact source support and finite $L^2$ norm give
absolute integrability and justify interchanging the integrals.
The same weighting convolves the finite-pulse spectral power and preserves
its integrated weight. Since the weight is even, the odd detuning response
cancels. The positive local slope found above therefore does not determine
whether phase diffusion increases or decreases mean work. Nor does a single
symmetric detuning pair settle this question.

For computation of this mean work, the lag factor $e^{-D(t-s)}$ can instead
be included in the auxiliary generator $A-DI$, evaluated for two genuine
source quadratures. With state coordinates $(q,p)$, both equations acquire
$-D$ terms; the work output remains $p+i\omega q$, although now
$\dot q=p-Dq$. These auxiliary states are not stochastic field trajectories,
and their endpoint energies are not ensemble energies. This reduction
concerns the specified work observable and does not insert physical friction
into the original field equations.

A Cauchy-distributed constant detuning and a Wiener phase process share the
relevant two-point statistics but not their trajectory laws. Their quadratic
response means agree; the auxiliary damped endpoint is a different object.
Ideal Wiener phase also has no finite ultraviolet cutoff. Thus this identity
is a finite-model interpretation, not a continuum certificate, a temperature
assignment, or a derivation of an environmental bath. Physical bandwidth and
coupling require separate specification before experimental use.
\paragraph{One finite-coherence comparison.}
We fixed $D=1$, corresponding to one pulse duration for the exponential
phase-correlation time, and compared it with $D=0$. The two genuine
quadratures of each auxiliary response were evolved on the same two
discretizations as the detuning test. The finer grid gives mean source
works $m_0=1.133373236\times10^{-4}$ and
$m_1=2.173314375\times10^{-4}$. Their difference,
$1.039941139\times10^{-4}$, is positive on both grids and exceeds the
predeclared diagnostic guard $6.63\times10^{-9}$. The undiffused adapter
reproduces the earlier work values, and the modified auxiliary balances
include $-2D$ times each quadratic form. Free and restoring-oscillator
controls test both the work output and the placement of damping in both
auxiliary coordinates.

This is an increase in mean work supplied by the prescribed external
source at fixed source $L^2$ budget. It is not a gain in energetic efficiency
or a demonstration of self-sustained oscillations. The guard combines
grid differences and implementation residuals; it is not a certified
error bound, and space and time resolution change together. The nominal
profile, finite boundary and linear-response limitations persist. Only
one nonzero diffusion constant is compared: no noise optimum, temperature,
nonlinear formation or lifetime conclusion is established.
The same source-output reduction with output $p$ gives an additional
observable. In the undamped physical finite model, the rotating quadratic
generator satisfies $\dot H=2\operatorname{Re}\langle p,S_\theta\rangle$.
For zero initial perturbation, its ensemble endpoint mean therefore equals
the half-sum of the two auxiliary integrals of this \emph{original} source
rate. It does not equal the auxiliary endpoint $H$. Extracting these
already stored integrals gives, on the finer grid,
$\mathbb E H(T)=2.09020168\times10^{-4}$ for $D=0$ and
$2.93069596\times10^{-4}$ for $D=1$. Both grids show this approximately
$40\%$ increase; this is a descriptive follow-up without a new certified
comparison bound or field integration.

This rotating quadratic generator need not be positive, and its global
value does not measure energy localized near the Q-ball. Full laboratory
endpoint energy additionally requires the physical ensemble charge,
including the anomalous state-exchange contribution. The latter is not
given by the exchange evaluated on the damped auxiliary states. Thus
neither the source-work increase nor this rotating-generator mean establishes
laboratory energy retention or improved confinement.

\paragraph{Correlated instantaneous frequency.}
A separate prescribed drive model replaces Wiener phase increments by
$\theta(t)=\theta_0+\int_\Delta^t\xi(s)\,ds$, with an independent uniform
initial angle and stationary zero-mean Ornstein--Uhlenbeck frequency
$\mathbb E\xi(t)\xi(s)=(D/\tau_c)e^{-|t-s|/\tau_c}$.
Its phase correlation at lag $u\ge0$ is
\[
 C(u)=\exp\{-D[u-\tau_c(1-e^{-u/\tau_c})]\}.
\]
This retains the pathwise source norm and fixes the long-time phase
diffusion constant $D$; the instantaneous frequency variance is
$D/\tau_c$. The Wiener kernel is recovered as $\tau_c\to0$ at fixed $D$.
Finite $\tau_c$ removes its linear short-lag cusp, but does not impose a
hard bandwidth cutoff or define a thermal bath. Pointwise larger phase
correlation need not give larger work against an oscillatory response.

\paragraph{Small correlation-time limit.}
For a fixed finite spatial operator and fixed $D\geq0$, let $\tau=\tau_c$ and write the exact mean work as
$m_\tau=\int_0^T W(u)C_\tau(u)\,du$, with real lag kernel $W$ and
$m_D=\int_0^T W(u)e^{-Du}\,du$. If $|W(u)|\leq M$, then
\[
 |m_\tau-e^{D\tau}m_D|\leq
 \frac{e^{D\tau}DM\tau^2}{1+D\tau},\qquad m'_+(0)=Dm_D.
\]
If $W$ is continuous at zero, the present source-work observable has
$W(0)=2B_{\mathrm{src}}$, where $B_{\mathrm{src}}=\int_0^T\|S(t)\|^2dt$ is the exact
nominal source budget. Consequently,
$m_\tau=m_D+Dm_D\tau+(D^2m_D/2-2DB_{\mathrm{src}})\tau^2+o(\tau^2)$.
These local asymptotics neither determine the comparison at $\tau=0.1$
nor imply global monotonicity, a thermal interpretation, or a continuum error bound.

For the fixed choice $D=1$, $\tau_c=1/10$, put $a=1/10$.
The mean source work has the signed representation
$m_{\mathrm{OU}}=e^a\sum_{n\ge0}(-a)^n m_{1+10n}/n!$, where $m_\lambda$
is the two-quadrature work obtained from the auxiliary generator
$A-\lambda I$. This is not a positive mixture of dissipative states.
The selected truncation retains $n=0,\ldots,8$. If
$B\ge2\int_{t>s}|K_{\mathrm{work}}(t,s)|\,dt\,ds$, its analytical work-tail
bound is
\[
 \epsilon_{\mathrm{tail}}\le
 B\,\frac{10}{9}\,\frac{(1/10)^9}{9!}\,\frac{1}{1-1/100}.
\]
The finite-model bound uses a positive reference energy norm,
a coefficient-envelope check and an exact-rational source-norm bound
for frozen nominal binary coefficients. This bounds the omitted series,
not floating exponential evaluation, signed-sum roundoff or RK integration.
Those errors and the combined space/time sensitivity require separate
assessment; auxiliary endpoint energies are not ensemble energies.
The first finite-correlation implementation passed its scalar controls but
failed the preselected modified quadratic-balance tolerance on the coarser
field grid. That implementation therefore yielded no accepted finite-correlation
field result, and its finer grid was not run. The analytical series-tail estimate
did not repair this numerical acceptance failure, and the original outcome
and thresholds remain unchanged. A separate calculation with a different
integrator and a prospectively specified temporal-sensitivity comparison is
reported below.

\paragraph{A separate finite-correlation pulse comparison.}
A separately specified two-stage Gauss calculation retained the prescribed
source, its normalization, $D=1$, $\tau_c=0.1$ and the signed nine-term expansion.
It split integration steps at the two known source-window joins and used
both genuine quadratures for every auxiliary rate, plus a zero-source
control. Three fixed integrations compared $(h,\Delta t)=(0.04,0.005)$,
$(0.04,0.0025)$ and $(0.02,0.0025)$ through $T=3$ at $R=20$.
All 108 recorded acceptance checks passed; this comparison was made against
the previously stored Wiener-source work on the corresponding spatial grid,
not against an auxiliary endpoint energy or a newly repeated Wiener run.

On the finer spatial grid, the truncated OU mean source work was
$2.1932622828\times10^{-4}$, compared with $2.1733143747\times10^{-4}$
for the stored Wiener result. The difference, $1.9947908110\times10^{-6}$
(approximately $0.918\%$), exceeded the predeclared diagnostic guard
$1.2658579115\times10^{-8}$ and had the same positive sign on the coarser grid.
The guard includes the analytical series-tail allowance and measured temporal,
spatial, adapter and accounting contributions; it is not a certified total-error
bound. The fixed coarse-grid time comparison does not directly measure fine-grid
temporal error. This is a bounded increase in mean source work for one prescribed
drive in the frozen finite-domain model, not evidence of retained energy,
heating, formation, improved efficiency or experimental realizability.

\begin{figure}[htbp]
\centering
\includegraphics[width=.98\linewidth]{mean-injected-work.pdf}
\caption{Cumulative mean source work for prescribed OU and Wiener phase
drives on the $h=0.02$ grid (left), and their difference (right).
The OU curve is the signed nine-term mean of the two genuine quadrature
works per auxiliary rate; the Wiener curve uses the previously stored
calculation. Both use integrated original source work, not auxiliary
endpoint or ensemble energy. Straight lines join 31 stored times, with no
additional field integration for the figure. The endpoint marker spans
the predeclared diagnostic guard; it is neither a confidence interval nor
a certified total-error bar. All quantities are in dimensionless model units.}
\label{fig:ou-source-work}
\end{figure}

\paragraph{One exterior-domain sensitivity check.}
A separate calculation at fixed $h=0.04$, $\Delta t=0.0025$ and $T=3$
extended $R=20$ to $R=30$, retaining byte-identical source values and
interior coefficients and comparing with the stored $R=20$ calculation.
Across the 31 shared sample times, the absolute-weighted majorant for the
change in the same-method OU-minus-$n=0$ source-work contrast was
$1.14\times10^{-19}$, below the predefined diagnostic threshold
$1.27\times10^{-9}$. This is agreement at floating-point roundoff scale,
not absolute accuracy at $10^{-19}$; common spatial, temporal and profile
errors can remain. The check does not establish an infinite-domain limit,
long-time behavior or control between samples, and neither replaces the
earlier Wiener baseline nor reduces the original comparison guard.

\paragraph{Signed contributions from source-time separations.}
A separate response-kernel calculation uses the fixed $h=0.04$, $R=20$
spatial operator and both genuine quadratures of each source profile.
After averaging the uniform initial phase, its mean source work is
$m[C]=\int_0^1 W(u)C(u)\,du$, where $u$ is the separation of two source
times, not the observation time. The kernel can have either sign, with
$W(0)=2B_{\mathrm{src}}$. On the finest of three lag meshes
($400$, $800$, $1600$ subdivisions), with independent overlap quadratures
of orders $16$ and $32$, the four prescribed quarter-interval contributions
to $\int_0^1 W(u)[C_{\mathrm{OU}}(u)-e^{-u}]\,du$ are, in units of $10^{-6}$,
\[
 (+23.13074873,\;-9.08245885,\;-11.01641674,\;-1.03681136).
\]
Their signed sum is $1.99506177879\times10^{-6}$, with summed heuristic
refinement envelopes $6.98\times10^{-14}$; these are not a certified
error bound or confidence interval. The shortest-lag contribution exceeds
the three negative later-lag contributions (Fig.~\ref{fig:signed-lag-kernel}).
Closure to the stored same-grid Gauss OU and Wiener-rate works passes the
predeclared diagnostic allowances. The separate shift of the Gauss
Wiener-rate work relative to the old RK Wiener baseline is
$-8.47524\times10^{-10}$; it is not assigned to any lag bin.
This calculation is not the $h=0.02$ comparison above and does not exactly
decompose its mixed-integrator result. A lag bin integrates pairs of source
times, not instantaneous positive or negative power. All separations are
at most one model time unit, below the nominal period $2\pi/\rho\simeq3.44$;
the response involves the full generator, not an isolated BIC or a long
ringdown. No stored-energy, formation or stabilization conclusion follows.

\begin{figure}[htbp]
\centering
\includegraphics[width=.98\linewidth]{signed-lag-kernel.pdf}
\caption{Stored finest-grid signed lag kernel $W(u)$ (left) and four fixed
integrals of $W(u)[C_{\mathrm{OU}}(u)-e^{-u}]$ (right), for the frozen $h=0.04$
model. Lines connect the stored lag nodes without smoothing. The caps and
reported net envelope are heuristic numerical sensitivity measures, too
small to resolve at this scale, not confidence intervals or certified
errors; no continuous kernel error band is asserted. The positive net
source-work difference retains all three negative contributions. These
are source-time-pair contributions, not a time history of instantaneous
power or retained energy.}
\label{fig:signed-lag-kernel}
\end{figure}


Scalar radiation pressure has been studied in a related
one-dimensional Q-ball model~\cite{Ciurla2024}; its rates are not imported
into our three-dimensional problem.

Several Q-balls of the same nonlinear field can influence one another,
but a sum of isolated stationary profiles is not automatically a stationary
multi-centre solution. Separation, relative phase and incident frequencies
must be specified. We have not derived a self-consistent environmental
noise spectrum, an interaction with every particle species, or a mechanism
that maintains these excitations indefinitely.

If an additional dynamical mediator is introduced, its spatially nonlocal
static response must be distinguished from temporal memory. Matching a
homogeneous effective potential does not generally match the nonlocal static
interaction. Such extensions are not included in Eq.~\eqref{eq:action}.

An additional open channel also changes the localization condition. For a
real-frequency mode with independent propagating asymptotic channels, every
propagating tail must vanish separately. A cancellation of the original
outgoing amplitude therefore does not establish an embedded bound state of
the extended system. Nor does a zero of a leading-order source overlap
establish this: the coupled background and mode can change at higher order.
These are conditions for a future coupled calculation, not results for an
extension of the model studied here.
